\phantomsection
\addcontentsline{toc}{section}{References}
\begin{thebibliography}{99}
\raggedright

\bibitem{sourcewiki}
Wikipedia contributors.
\emph{Propositional logic}, and the immediate mathematical links listed in
Appendix~\ref{sec:source-map}. Inspected 22 November 2026.
\href{https://en.wikipedia.org/wiki/Propositional_logic}{Starting article}.
Used to specify the requested subject coverage; the mathematical proofs are
given in the report.

\bibitem{proveit}
Vladimir Reshetnikov and contributors.
\emph{ProveIt}, especially \texttt{Logic/PrimitiveRecursiveArithmetic}.
Repository snapshot
\texttt{803d951d15834da5571e26d1f39988d921b29760}, inspected
22 November 2026.
\href{https://github.com/VladimirReshetnikov/ProveIt/tree/803d951d15834da5571e26d1f39988d921b29760/Logic/PrimitiveRecursiveArithmetic}{Pinned arithmetic directory}.
The inspected declarations and their intended integration roles are listed
in Section~\ref{sec:lean}.

\bibitem{kripke1965}
Saul A. Kripke.
``Semantical Analysis of Intuitionistic Logic I.''
In J. N. Crossley and M. A. E. Dummett, editors,
\emph{Formal Systems and Recursive Functions},
Studies in Logic and the Foundations of Mathematics, volume 40,
North-Holland, 1965, pp.~92--130.
\href{https://doi.org/10.1016/S0049-237X(08)71685-9}{DOI:
10.1016/S0049-237X(08)71685-9}.

\bibitem{buss-pudlak2001}
Samuel R. Buss and Pavel Pudl\'ak.
``On the Computational Content of Intuitionistic Propositional Proofs.''
\emph{Annals of Pure and Applied Logic}, volume 109, 2001.
\href{https://mathweb.ucsd.edu/~sbuss/ResearchWeb/intuitionisticDP2/paper.pdf}{Author-hosted paper};
\href{https://mathweb.ucsd.edu/~sbuss/ResearchWeb/intuitionisticDP2/index.html}{author's publication page}.
Section 4 of the paper treats disjunction extraction under Harrop
assumptions. The report proves its own weaker, primitive recursive bound.

\bibitem{burgess1981}
John P. Burgess.
``The Completeness of Intuitionistic Propositional Calculus for its Intended
Interpretation.''
\emph{Notre Dame Journal of Formal Logic} 22(1), 1981, pp.~17--28.
\href{https://doi.org/10.1305/ndjfl/1093883336}{DOI:
10.1305/ndjfl/1093883336};
\href{https://www.cs.cornell.edu/courses/cs4860/2019fa/lectures/Completeness_of_intuitionistic_propositional_logic.pdf}{Cornell-hosted copy}.
Cited to distinguish finite Kripke completeness from completeness for an
intended interpretation using additional principles.

\bibitem{cook1971}
Stephen A. Cook.
``The Complexity of Theorem-Proving Procedures.''
\emph{Proceedings of the Third Annual ACM Symposium on Theory of Computing},
1971, pp.~151--158.
\href{https://doi.org/10.1145/800157.805047}{DOI:
10.1145/800157.805047};
\href{https://www.cs.toronto.edu/~sacook/homepage/1971.pdf}{Author-hosted original paper}.

\bibitem{DP1960}
Martin Davis and Hilary Putnam.
``A Computing Procedure for Quantification Theory.''
\emph{Journal of the ACM} 7(3), 1960, pp.~201--215.
\href{https://doi.org/10.1145/321033.321034}{DOI:
10.1145/321033.321034};
\href{https://web.stanford.edu/class/cs357/DP60.pdf}{Stanford-hosted original paper}.
Only the finite propositional elimination construction is developed in
the present report; no first-order completeness theorem is imported.

\bibitem{Tseitin1968}
G. S. Tseitin.
``On the Complexity of Derivation in Propositional Calculus.''
Original Russian publication in \emph{Zapiski Nauchnykh Seminarov LOMI},
volume 8, 1968, pp.~234--259.
English translation in A. O. Slisenko, editor,
\emph{Studies in Constructive Mathematics and Mathematical Logic}, Part II,
Consultants Bureau, New York, 1970, pp.~115--125.
\href{https://mathsoc.spb.ru/pers/tseytin/bib.html}{Publication bibliography};
\href{https://yurichev.com/mirrors/SAT/Tseitin/Tseitin1968_en.pdf}{English reprint}.
The gate-clause equivalences and unique-extension property used here are
proved directly in Section~\ref{ca-solvers}.

\end{thebibliography}
